\textbf{Figure 2. Fine-mapping and liver-eQTL colocalization narrow liver-trait GWAS signals to candidate genes for testing in other measurements.}

Genetic counts, overlap estimates, accessibility context and Unified Gene Catalog classes reflect the allele-aware, positive-semidefinite (PSD) LD-projected eQTL refit of the 90 genes with duplicate variant identifiers (\texttt{EQTL-DEDUP-PSD-20261118}: 84 kept positive, 4 lost [CDK5RAP3, KANSL1-AS1, LRRC37A, THBS3], 2 untestable [CTSS, MPRIP]; 430 primary multi-signal SuSiE genes and 990 union genes; 5/5 unaffected controls pass parity at \textbar $\Delta$PP.H4\textbar{} $\le$ 2.55 $\times$ 10$^{-4}$).


\textbf{(A)} Liver-trait GWAS to colocalized genes. Left to right: 35 Tier-1/2 GWAS strata by ancestry (EUR 17, AFR 6, EAS 6, AMR 3, SAS 3) and trait; 360 trait-specific fine-mapped loci (265 unique) by trait; and 990 genes colocalized with liver eQTL after the PSD eQTL refit (PP.H4 \textgreater{} 0.5), split into 430 multi-signal (SuSiE) genes and 560 single-signal-only (ABF) genes and shaded by the band of each gene's largest PP.H4 (\textgreater{} 0.9, 0.8--0.9, 0.5--0.8: 239, 89, 102 multi-signal; 67, 89, 404 single-signal only). Three genes drawn as single-signal only (CKB, GSTM2, LILRB3) could not be tested by the multi-signal method; they are not multi-signal negatives. Loci are merged credible-set windows within each trait, so a locus tagged by several enzymes recurs across traits. MVP European strata carry colocalized genes but no credible-set loci and are drawn as faint trait-to-gene ribbons.

\textbf{(B)} Largest multi-signal PP.H4 per gene against SuSiE PIP of the same lead variant in the same GWAS (n = 865 genes). Blue, noncoding lead (309); magenta, coding or splice lead (16; 5 in the colocalized gene, 11 in another gene); light gray, 540 genes outside the display filter (PP.H4 $\le$ 0.05 and PIP \textless{} 0.99). Dashed line, PP.H4 = 0.5. Genes whose largest posterior comes from the unfine-mapped MVP European strata have no PIP and are not drawn; absence is not a negative call. Outlined and labeled, the three recurring example genes (Figure S2AC) that have a PIP in their best GWAS: PPM1K (PIP 0.002, PP.H4 0.95, GGT), HKDC1 (0.16, 0.99, AST) and CFHR5 (0.003, 0.73, ALT). HSD17B13 and RORA are labeled for reference (RORA colocalization rests on LD with a nearby locus; Methods).

\textbf{(C)} Consequence class of one lead COLOC variant per gene (990 genes, 927 distinct variants; lead = highest PP.H4, SuSiE or ABF), first matching hg19 knownGene class in the order coding or splice, UTR, promoter, intron, intergenic, with intron and intergenic from direct transcript overlap. Coding includes synonymous; 13 of the 34 coding leads fall in the colocalized gene and 21 in another gene. Counts: coding 34, UTR 70, promoter 104, intron 522, intergenic 260. A gene-level summary, not the posterior composition of credible sets.

\textbf{(D)} Ancestry of the GWAS supporting each of the 430 multi-signal genes (SuSiE PP.H4 \textgreater{} 0.5), counted per gene, not locus. Area is proportional to genes: EUR only, 334 (78\%); shared EUR and non-EUR, 64 (15\%); non-EUR only with lead p $\ge$ 5 $\times$ 10$^{-8}$, 23 (5\%); non-EUR only with lead p \textless{} 5 $\times$ 10$^{-8}$, 9 (2\%; 8 EAS, 1 AFR). Of the 334 EUR-only genes, 121 have a finite recorded non-EUR SuSiE output at PP.H4 $\le$ 0.5 and 213 have none; absent output is not a negative result. The eQTL panel is European, so non-EUR GWAS are not ancestry-matched.

\textbf{(E)} All 13 genes with PP.H4 \textgreater{} 0.5 (SuSiE or ABF) in four of five ancestry groups. Circle area, maximum PP.H4 over each group's studies. Dark blue, SuSiE; light blue, ABF only; gray, neither above 0.5. Asterisk, four groups by SuSiE alone. These partitions do not establish biological differences between ancestry groups.

\textbf{(F)} GWAS trait that gives each of the 430 candidates its largest SuSiE-coloc PP.H4, each gene counted once: AST 155, ALT 135, GGT 127, NAFLD 8, PDFF 4, NASH 1 (417 liver-enzyme and 13 direct MASLD/liver-fat drivers). The 13 direct-trait genes are ACSL5, ADH1C, CENPQ, ERCC2, IL18R1, PNPLA6, RABGEF1, RANBP17, SAMM50, SHMT1, UGT1A1, UGT1A8 and ZBTB41. The label records where a gene's strongest evidence comes from; it does not test a disease effect. Per-locus MASLD liability of enzyme-colocalized loci is in Figure S2Z--S2AB.

\textbf{(G)} Of 28,233 Gene Catalog genes, 18,724 are testable for liver-eQTL colocalization (of the other 9,509, 9,312 are underpowered in the eQTL source, 195 have insufficient shared posterior and 2 became untestable in the PSD eQTL refit) and 430 have multi-signal SuSiE-coloc PP.H4 \textgreater{} 0.5 in at least one of 35 Tier-1/2 GWAS strata. Measurements in which the 430 candidates can be tested next. Light bar, measured or detected; dark, 26 genes that change in bulk MASLD liver (canonical bulk call; Figure 3); gray, not measured. Bulk liver RNA, 420 in the five-cohort analysis; single-cell RNA, 370 detected in at least one of 9 donor-collapsed lineages (104 donors, 7 datasets); liver proteome, 241 detected in the PXD051911 liver DIA proteome (58 participants); CosMx, 29 targeted by the Govaere CosMx panel (one through a shared probe; panel design, not detection); plasma proteome, 19 detected in the PXD051911 plasma proteome (143 participants at baseline). Coverage does not imply a result.
